% !TeX program = xelatex
\documentclass[12pt,a4paper]{extarticle}

% 1. Подключение общей преамбулы (пакеты)
\input{../common/preamble}

% 2. Определение переменных для конкретного документа
\newcommand{\docname}{Текст программы}
\newcommand{\doccode}{\hseDocCodeBase\ 12 01-1}
\newcommand{\doccodelu}{\hseDocCodeBase\ 12 01-1-ЛУ}

% 3. Подключение общих команд (проект, студент, руководитель)
\input{../common/commands}

% 4. Подключение стилей (колонтитулы, заголовки)
\input{../common/styles}

\begin{document}
\sloppy

% 5. Генерация титульных листов по шаблону
\input{../common/title_template}


% ============================================
% АННОТАЦИЯ
% ============================================

\section*{АННОТАЦИЯ}
\addcontentsline{toc}{section}{АННОТАЦИЯ}

% Подсказка: перечислите, какие сведения о программе приведены в документе
% (карта репозитория, назначение сервисов, библиотеки, контракты и т. п.).
В настоящем документе «\projectname. Текст программы» приведены сведения о составе, иерархической структуре и используемом программном обеспечении разрабатываемой системы. Документ содержит карту репозитория, назначение реализуемых компонентов, описание состава общих библиотек и описание инфраструктурного слоя.

% Подсказка: фраза о закрытом репозитории опциональна — удалите её,
% если ваш репозиторий публичный.
Полный исходный код программы расположен в удаленном Git-репозитории. \hseOptional{Репозиторий является закрытым; доступ к нему предоставляется по запросу для проверки состава и реализации программы.}

\clearpage

% ============================================
% СОДЕРЖАНИЕ
% ============================================
\hypersetup{linkcolor=black}
\tableofcontents
\hypersetup{linkcolor=blue}
\thispagestyle{fancy}
\clearpage

% ============================================
% 1. ВВЕДЕНИЕ
% ============================================
\section{ВВЕДЕНИЕ}

\subsection{Наименование программы}

Наименование программы~--- «\projectname».

% Английское и краткое наименования подставляются автоматически из
% настроек проекта (поля «Название работы на английском» и «Краткое
% наименование программы») — здесь править не нужно.
Наименование программы на английском языке~--- «\hseProjectNameEn».

Краткое наименование программы~--- «\hseProjectNameShort».

\subsection{Краткая характеристика области применения программы}

% Подсказка: в 2-3 предложениях опишите назначение программы и область
% её применения; укажите архитектурный стиль (монолит, микросервисы,
% монорепозиторий) и ключевые технологии.
Программа «\projectname» является \hseFill{краткая характеристика области применения}. Она реализует механизмы \hseFill{перечень ключевых функций программы}. Система спроектирована по \hseFill{архитектурный стиль} архитектуре и использует \hseFill{ключевые технологии и протоколы}.

\clearpage

% ============================================
% 2. ТЕКСТ ПРОГРАММЫ
% ============================================
\section{ТЕКСТ ПРОГРАММЫ}

% Подсказка: укажите, что полный текст программы расположен в репозитории,
% и при необходимости приведите адрес и условия доступа.
% Подсказка: фраза о закрытом репозитории опциональна — удалите её,
% если ваш репозиторий публичный.
Полный текст программы расположен в удаленном Git-репозитории по следующему адресу: \hseFill{адрес Git-репозитория}. \hseOptional{Репозиторий является закрытым; доступ к нему предоставляется по запросу для проверки состава и реализации программы.}

\subsection{Структура репозитория}

% Подсказка: опишите верхнеуровневую структуру каталогов проекта.
% Перечислите назначение основных директорий своего репозитория.
\begin{hseExample}
Программный комплекс организован по принципу \hseFill{принцип организации репозитория}. Верхнеуровневая структура директорий:

\begin{itemize}[leftmargin=1.25cm]
    \item \texttt{src/} --- исходный код приложения: сервисы, модули и фоновые обработчики;
    \item \texttt{libs/} --- общие библиотеки, клиентские SDK и компоненты среды выполнения;
    \item \texttt{infra/} --- код управления инфраструктурой, конфигурации и утилиты развертывания;
    \item \texttt{docs/} --- исходные тексты технической, продуктовой и эксплуатационной документации;
    \item \texttt{scripts/} --- вспомогательные скрипты разработки, сборки и генерации артефактов;
    \item \texttt{tests/} --- интеграционные и системные тестовые материалы.
\end{itemize}
\end{hseExample}

\subsection{Описание модулей (сервисов)}

% Подсказка: для каждого ключевого модуля (сервиса) приведите его
% назначение. Заполните таблицу ниже своими данными.
\begin{hseExample}
Модули (сервисы) программы разделены на домены: \hseFill{перечень доменов или групп модулей}. Назначение ключевых модулей приведено в таблице~\ref{tab:modules}.

% Таблица набрана через \captionof (не \begin{table}[H]): плавающие
% окружения внутри жёлтого блока-образца недопустимы.
\begin{center}
\small
\captionof{table}{Назначение модулей программы}
\label{tab:modules}
\begin{tabularx}{\textwidth}{|>{\raggedright\arraybackslash}p{0.32\textwidth}|X|}
\hline
\textbf{Модуль} & \textbf{Назначение} \\
\hline
\texttt{\hseFill{имя модуля}} & \hseFill{краткое назначение модуля} \\
\hline
\texttt{\hseFill{имя модуля}} & \hseFill{краткое назначение модуля} \\
\hline
\end{tabularx}
\end{center}
\end{hseExample}

\subsection{Общие библиотеки}

% Подсказка: опишите используемые общие (внутренние) библиотеки и их
% назначение. Если общих библиотек нет, раздел можно удалить.
\begin{hseExample}
Для обеспечения единообразия разработки и переиспользования кода используется набор внутренних библиотек, сгруппированных по назначению.

\begin{center}
\small
\captionof{table}{Назначение групп общих библиотек}
\label{tab:libs}
\begin{tabularx}{\textwidth}{|>{\raggedright\arraybackslash}p{0.32\textwidth}|X|}
\hline
\textbf{Группа} & \textbf{Состав и назначение} \\
\hline
\texttt{\hseFill{группа библиотек}} & \hseFill{состав и назначение группы} \\
\hline
\texttt{\hseFill{группа библиотек}} & \hseFill{состав и назначение группы} \\
\hline
\end{tabularx}
\end{center}
\end{hseExample}

\subsection{API-контракты и клиентское взаимодействие}

% Подсказка: опишите способ внешнего и межсервисного взаимодействия
% (REST, gRPC, форматы данных, расположение схем контрактов).
\begin{hseExample}
Внешнее взаимодействие с программой описывается через \hseFill{протоколы внешнего API}. Точки доступа принимают и возвращают данные в формате \hseFill{формат данных}, используемые клиентскими системами. Межсервисное взаимодействие строится на \hseFill{механизм межмодульного взаимодействия}: исходные файлы схем контрактов хранятся в каталоге \texttt{\hseFill{каталог со схемами}}, а клиентские и серверные заглушки генерируются \hseFill{средство генерации кода}.
\end{hseExample}

\subsection{Хранение данных и событийное взаимодействие}

% Подсказка: укажите используемые СУБД, средства кэширования и обмена
% сообщениями, а также применяемые паттерны хранения и доставки данных.
\begin{hseExample}
Серверные компоненты реализованы преимущественно на \hseFill{укажите языки и фреймворки}. Для хранения бизнес-данных применяется \hseFill{используемая СУБД}, для кэширования и краткоживущего состояния~--- \hseFill{средство кэширования}, для асинхронного обмена событиями~--- \hseFill{брокер сообщений}.

% Подсказка: опишите, как обеспечивается надежная доставка событий и
% согласованность данных (например, Transactional Outbox/Inbox);
% если событийного обмена нет, абзац можно удалить.
Надежная доставка событий обеспечивается \hseFill{применяемые паттерны доставки событий}: доменное изменение и запись события выполняются в одной транзакции, после чего фоновый процесс публикует событие в \hseFill{брокер сообщений}.
\end{hseExample}

\subsection{Инфраструктура и развёртывание}

% Подсказка: кратко опишите инфраструктурный слой: контейнеризацию,
% средства оркестрации, автоматизацию сборки и развёртывания.
\begin{hseExample}
Инфраструктурный слой обеспечивает автоматизацию развертывания и эксплуатации:

\begin{itemize}[leftmargin=1.25cm]
    \item \textbf{\hseFill{каталог средств оркестрации}} --- манифесты, конфигурации и эксплуатационные скрипты для \hseFill{компоненты инфраструктуры};
    \item \textbf{scripts/} --- CI/CD и эксплуатационные скрипты;
    \item \textbf{docker-compose.yml} --- локальное инфраструктурное окружение для разработки и отладки.
\end{itemize}
\end{hseExample}

\subsection{Контракты и документация}

% Подсказка: опишите, где хранятся схемы контрактов и проектная
% документация; если раздел не актуален, его можно удалить.
\begin{hseExample}
\begin{itemize}[leftmargin=1.25cm]
    \item \textbf{\hseFill{каталог схем контрактов}} --- иерархическая структура файлов, определяющая интерфейсы \hseFill{перечень основных интерфейсов};
    \item \textbf{docs/} --- структурированная база знаний, разделенная по доменам: \hseFill{перечень доменов документации};
    \item \textbf{README.md} и сервисные \textbf{README.md} --- инструкции по запуску, разработке, составу API и эксплуатационным особенностям отдельных компонентов.
\end{itemize}
\end{hseExample}

\clearpage
\pagestyle{appendix}

% ============================================
% 3. СПИСОК ИСПОЛЬЗОВАННЫХ ИСТОЧНИКОВ
% ============================================
\section*{СПИСОК ИСПОЛЬЗОВАННЫХ ИСТОЧНИКОВ}
\addcontentsline{toc}{section}{СПИСОК ИСПОЛЬЗОВАННЫХ ИСТОЧНИКОВ}

% Подсказка: приведите перечень использованных источников. Ниже оставлены
% обязательные нормативные документы ЕСПД и один пример источника на ПО —
% замените или дополните своими.
\renewcommand{\refname}{}
\begin{thebibliography}{99}
\vspace{-1cm}

\bibitem{python_docs}
Python Documentation [Электронный ресурс]. Режим доступа: \url{https://docs.python.org/3/}, свободный. (дата обращения: 01.01.2026).

\bibitem{gost19101}
ГОСТ 19.101-77: Виды программ и~программных документов. // Единая система программной документации.~--- М.: ИПК Издательство стандартов, 2001.

\bibitem{gost19105}
ГОСТ 19.105-78: Общие требования к~программным документам. // Единая система программной документации.~--- М.: ИПК Издательство стандартов, 2001.

\bibitem{gost19401}
ГОСТ 19.401-78: Текст программы. Требования к~содержанию и~оформлению. // Единая система программной документации.~--- М.: ИПК Издательство стандартов, 2001.

\end{thebibliography}



\clearpage

% ============================================
% ЛИСТ РЕГИСТРАЦИИ ИЗМЕНЕНИЙ
% ============================================
\input{../common/change_log}

\end{document}
